\section{Conclusion and outlook}
\label{sec:conclusion}

\subsection{Technical synthesis}

The starting observation was that an optimiser can succeed while a registration
fails. The investigation traced this to one structural fact: the scan supplies
surface geometry, whereas anatomical correspondence is an intended consequence of
the recovered parameters and is constrained only through the model, its priors and
independent verification. Four distinctions summarise the
outcome: representability is not recoverability; regional identity is not
within-region location; plausibility is not correspondence; and fit quality is not
measurement validity. Each was the concrete reason that a plausible intervention
failed to deliver what it was built for.

Technically, the investigation extended the pipeline across mesh conditioning,
anatomical priors, learned point-cloud initialisation, sampling-aware and robust
optimisation, learned correspondence and hierarchical fitting, while introducing
synthetic and expert-grounded instruments to evaluate these mechanisms
independently. Taken together, the experiments indicate that the dominant
difficulty is not simply surface alignment but recovery of a unique anatomical
explanation from an observation that lacks explicit identity information. A useful
parametric representation is defined not only by what it can represent but by what
anatomical meaning survives its recovery from data: a reconstruction is not yet a
phenotype.

\subsection{Research directions}

Each direction follows from a diagnosed limitation (Table~\ref{tab:roadmap}), and
none was implemented.

\begin{table}[!htb]
\centering
\caption{Diagnosed limitations and the research directions they motivate.}
\label{tab:roadmap}
\begin{tabular}{@{}L{0.40\linewidth}L{0.52\linewidth}@{}}
\toprule
\textbf{Diagnosed limitation} & \textbf{Next technical direction}\\
\midrule
Local votes lack global consistency & Functional maps~\cite{fmaps}, optimal
 transport~\cite{ot}, globally coupled assignment\\
Joint position is not observed & Skeleton-conditioned correspondence, a learned 3D
 joint predictor\\
Free deformation can absorb pose & Structured, diffeomorphic deformation\\
Synthetic learning transfers unevenly & Scan-realistic corruption and domain
 randomisation\\
Part identity does not give precise location & A continuous anatomical canonical
 field with uncertainty, in the spirit of probabilistic implicit
 embeddings~\cite{liu2020}\\
\bottomrule
\end{tabular}
\end{table}

The most ambitious direction is an anatomical canonical field that predicts, for
every observed point, its part, joint-relative coordinates, laterality and a
confidence, so that correspondence becomes the mapping of a scan into a canonical
anatomical coordinate system: exactly the information the diagnosed objective lacks.

\subsection{Broader significance}
\label{sec:significance}

The ant is a useful test case because it forces diversity, articulation, difficult geometry
and correspondence to be handled at once (Figure~\ref{fig:corpus}): genera differ in body
plan and posture, the structures of interest are thin and articulated, and the data are
imperfect. The aim, however, is broader than one good ant model. The longer-term objective
is to use what is learned here to build purpose-built articulated models for the body plans
that a scientific question actually requires, spanning other animals within a single
construction and fitting pipeline and not one predefined model family, as the mouse and
multi-species insect models already built with the same framework~\cite{smilify} begin to
show. It extends equally to human-centred applications, where generic parametric models can
impose restrictions that are inappropriate to the question being asked.

\begin{figure}[H]
\centering
\includegraphics[width=0.78\linewidth]{figures/fig_corpus_gallery.png}
\caption[Diversity of the curated corpus]{Eight genera from the curated corpus of clean
scans, rendered under identical conditions: the shared representation has to accommodate
differences in body plan, proportions, posture and articulation.}
\label{fig:corpus}
\end{figure}

The value of the work is therefore not only the fitted mesh. It is an understanding of how
to construct the representation, how to fit it, where the fitting can fail, and how to
establish what its parameters mean. The question that started the project was how to bring
many different meshes into one shared representation; the harder question it leads to is
what that representation can be trusted to mean. Once that is established, the value extends
beyond reconstruction: it becomes a way of making morphology measurable.
